\documentclass[11pt,a4paper]{article}

\usepackage[margin=1in]{geometry}
\usepackage{amsmath,amssymb,amsfonts}
\usepackage{mathtools}
\usepackage{lmodern}
\usepackage[T1]{fontenc}
\usepackage[utf8]{inputenc}
\usepackage{textcomp}
\usepackage{microtype}
\usepackage{parskip}
\usepackage{enumitem}
\usepackage{array}
\usepackage{booktabs}
\usepackage{longtable}
\usepackage{tabularx}
\usepackage{xcolor}
\usepackage{hyperref}
\hypersetup{hidelinks,
  pdftitle={Substrate-Level Resolution of the Black Hole Information Paradox},
  pdfauthor={Allen Proxmire}}

\setlength{\parindent}{0pt}
\setlength{\parskip}{6pt plus 2pt minus 1pt}

% Centered title block
\title{\vspace{-1cm}\Large\bfseries Substrate-Level Resolution of the Black Hole Information Paradox\\[6pt]
\large\mdseries Integrating Event Density Architecture with Hawking-Radiation Derivations}
\author{Allen Proxmire}
\date{May 2026}

\renewcommand{\abstractname}{\large Abstract}

\begin{document}
\maketitle
\thispagestyle{empty}

\vspace{-0.5cm}

\section*{Abstract}
The black-hole information paradox has resisted resolution within standard semiclassical and quantum-gravitational frameworks for nearly fifty years. Proposed resolutions --- black-hole complementarity, firewalls, ER=EPR, soft hair, replica wormholes and the island formula --- each address specific subproblems but each rests on structural assumptions (global unitarity, geometric horizon as primitive boundary, monogamy of entanglement enforced at that boundary, global Cauchy-data evolution) that the Event Density (ED) framework does not impose at the substrate level. This paper integrates the recently-closed Arc Hawking results (substrate-level derivation of the Hawking spectrum at $T_H = \kappa/(2\pi)$, V5 finite-memory cutoff at the Planck scale, Page rate $\dot M = -\alpha_{\mathrm{Page}}/M^2$, Page curve via the bandwidth-budget min-bound, \textbf{higher-order resummation establishing a stable Planck-mass remnant FORCED at substrate-structural level (H-8), and a cosmological relic-matter abundance calculation (H-9)}) with the previously established ED black-hole architecture. We show that the standard paradoxes are \emph{not generated} in the ED framework rather than \emph{resolved within it}. The substrate-level V5 cutoff resolves the trans-Planckian problem; the entanglement-straddling mechanism replaces the pair-creation picture; the bandwidth-budget min-bound replaces the island-formula entanglement-entropy bookkeeping; and substrate-discreteness eliminates the smooth-horizon idealization on which the standard paradoxes depend. \textbf{The framework predicts a structural Planck-mass remnant population from PBH evaporation with present-day relic-matter fraction $\Omega_{\mathrm{relic}}$ scenario-dependent on the empirical PBH formation history.} \textbf{Critical framing note:} ED's substrate-gravity (separate galactic\_dynamics walkthrough) explains galactic phenomenology --- flat rotation curves, slope-4 BTFR, transition acceleration $a_0 = cH_0/(2\pi)$ --- \emph{without} invoking dark matter. The Planck-mass remnant population is therefore a \emph{separate structural cosmological prediction} about cosmic relic-matter content, not a dark-matter candidate. We articulate ED's position as a \emph{regulated completion} of semiclassical Hawking with \textbf{definite Scenario C late-time content (Planck-mass remnant FORCED)}, distinguish it explicitly from each major proposal in the active BH-information literature, and identify falsifiable predictions accessible in current-generation analog Hawking experiments, primordial-BH evaporation observations, and high-energy photon dispersion measurements.

\textbf{Keywords:} black hole information paradox, Hawking radiation, Event Density, substrate ontology, Page curve, firewalls, ER=EPR, island formula, regulated completion, trans-Planckian problem.

\begin{center}\rule{0.5\textwidth}{0.4pt}\end{center}

\section{Introduction}

In 1975, Hawking demonstrated that black holes radiate thermally at temperature $T_H = \kappa/(2\pi)$, where $\kappa$ is the surface gravity of the horizon [1]. The result is one of the most striking convergences of general relativity and quantum mechanics: it ties together black-hole thermodynamics, Bekenstein's earlier proposal that black-hole entropy is proportional to horizon area [2], and the Bogoliubov machinery of quantum field theory in curved spacetime. Hawking's calculation was widely accepted as mathematically correct.

In 1976, Hawking himself articulated the structural problem the calculation generated [3]: if the radiation is purely thermal, the information content of any matter that fell into the black hole is destroyed when the black hole evaporates. Unitary evolution of a pure quantum state cannot produce a thermal mixed state, and this represents an apparent contradiction between general relativity (predicting evaporation to thermal radiation) and quantum mechanics (predicting unitary evolution preserving information).

The intervening five decades have generated an enormous theoretical literature, with proposed resolutions including:
\begin{itemize}[noitemsep]
  \item Black-hole complementarity [4]
  \item Firewalls [5]
  \item ER=EPR [6]
  \item Soft hair on black holes [7]
  \item The island formula and replica wormholes [8, 9]
\end{itemize}

Each of these proposals addresses subproblems within the standard framework. Each rests on structural commitments that constrain the available solution space in specific ways. None has achieved consensus acceptance.

The Event Density (ED) framework provides a structurally different approach. The framework is built on a substrate ontology of discrete micro-events on chains, with bandwidth, polarity, locality, finite-width kernels, and irreversible commitment-events as primitive objects. Within this ontology, the assumptions that generate the standard paradoxes --- global unitarity, geometric horizon as primitive boundary, monogamy at that boundary, global Cauchy-data evolution --- are not imposed. The phrase that captures the framework's stance: \emph{black holes do not have paradoxes because black holes do not have smooth edges, because the substrate is discrete}.

This paper integrates the recently-closed Arc Hawking machinery [10] with the ED black-hole architecture [11] to articulate the substrate-level resolution. The argument proceeds in three steps. First, we identify the structural assumptions in standard treatments that generate the paradoxes (\S2). Second, we show how the ED substrate ontology fails to make these assumptions and produces, via Arc Hawking, a substrate-level derivation of Hawking radiation that reproduces the empirically validated semiclassical content while resolving structural questions (\S\S3--8). Third, we compare the ED resolution explicitly with each major proposal in the active literature (\S\S9--10), articulating where ED agrees, where it differs, and what specific predictions distinguish it (\S\S11--13).

The paper's principal claim is that ED is best understood as a \emph{regulated completion} of standard semiclassical Hawking. The framework includes all of semiclassical Hawking at leading order, provides substrate-level UV regulation (V5 finite-memory cutoff at the Planck scale) that resolves the trans-Planckian problem, supplies a structurally economical account of information transfer (entanglement-straddling rather than pair creation), and produces conditional new content at extreme scales (Planck-mass remnant scenario contingent on higher-order analysis). The framework does not replace semiclassical Hawking; it grounds it.

\section{The Standard Paradoxes and Their Generating Assumptions}

\subsection{The original information-loss paradox}

Hawking's 1976 argument [3] runs through three steps:

\begin{enumerate}[label=(\roman*),noitemsep]
  \item Quantum mechanics requires unitary evolution of pure states.
  \item A black hole formed from a pure quantum state evaporates into Hawking radiation.
  \item The radiation, calculated in the Bogoliubov framework, is thermal --- a mixed state.
\end{enumerate}

If (i)--(iii) all hold, the apparent transition from pure to mixed state during BH evaporation violates unitarity and constitutes information destruction. The paradox is the inconsistency between quantum-mechanical unitarity (i) and the empirically-derived thermal character of Hawking radiation (iii) when combined with the macroscopic process of BH evaporation (ii).

The structural assumptions on which the paradox depends:

\textbf{(A1) Global unitarity} of the full quantum state on a single Hilbert space throughout evaporation.

\textbf{(A2) Geometric horizon} as a sharp boundary defining inside/outside.

\textbf{(A3) Information localization at horizon-crossing time} --- information falling in is ``stored'' inside the BH, and must somehow re-emerge through Hawking radiation.

\textbf{(A4) Pure-thermal Hawking spectrum} --- radiation is thermal in form, with no correlations between Hawking quanta that could carry information.

\subsection{Firewalls (AMPS 2012)}

Almheiri, Marolf, Polchinski, and Sully [5] showed that three reasonable assumptions --- (i) Hawking radiation is purely random thermal, (ii) the black-hole horizon is benign for infalling observers (equivalence principle), and (iii) the BH information is preserved unitarily --- are jointly inconsistent for sufficiently old black holes. Specifically, after the Page time, monogamy of entanglement requires that the early radiation be maximally entangled with the late radiation; this contradicts the requirement that infalling observers see the late-radiation modes maximally entangled with their interior partners (smooth horizon).

The firewall paradox replaces the original information paradox: rather than asking ``is information lost?'', it asks ``what does an infalling observer encounter at the horizon?'' If the equivalence principle holds, the observer sees nothing unusual; if monogamy is enforced, the observer encounters a wall of high-energy radiation.

The structural assumptions on which firewalls depend:

\textbf{(A5) Monogamy of entanglement enforced at the geometric horizon} --- a single mode cannot be maximally entangled with two systems simultaneously, so old-radiation/new-radiation entanglement contradicts new-radiation/interior-partner entanglement.

\textbf{(A6) Effective field theory smoothness across horizon} for an infalling observer --- the equivalence principle's empirical content.

The combination of (A1) + (A5) + (A6) is what produces the firewall.

\subsection{ER=EPR (Maldacena-Susskind 2013)}

Maldacena and Susskind [6] proposed that every entangled pair is connected by a (typically non-traversable) Einstein-Rosen bridge --- a wormhole. The proposal addresses firewalls by allowing the interior modes to be \emph{literally identified} with the exterior radiation modes via the wormhole geometry, dissolving the apparent monogamy violation: the modes are the same modes, viewed from different ends of the wormhole.

The structural assumptions on which ER=EPR depends:

\textbf{(A7) Wormhole topology} as physically realized for entangled pairs.

\textbf{(A8) Identification of interior and exterior modes} through the wormhole geometry.

\textbf{(A9) Geometric realization of entanglement} as connected geometry rather than purely correlational structure.

ER=EPR is a structurally elegant proposal but rests on the assumption that wormholes are physically realized. This is a strong ontological commitment for which independent evidence is unavailable.

\subsection{Soft hair (Hawking-Perry-Strominger 2016)}

Hawking, Perry, and Strominger [7] proposed that black holes carry ``soft hair'' --- supertranslation and superrotation charges associated with infinite-dimensional asymptotic symmetry groups. The soft-hair charges are gauge-invariant boundary data that distinguish black-hole microstates without modifying the bulk geometry. They proposed that information falling into the BH is encoded in soft-hair imprints on the horizon and asymptotic infinity, allowing information to escape via correlations with soft modes during evaporation.

The structural assumptions on which soft hair depends:

\textbf{(A10) Asymptotic symmetry algebra} as physically meaningful (BMS symmetries at null infinity, supertranslations at horizon).

\textbf{(A11) Soft-mode encoding} of bulk information in boundary degrees of freedom.

\textbf{(A12) Geometric horizon} retains its primitive role as a boundary where soft charges are defined.

Soft hair has been extensively developed but has not produced a complete account of how information actually escapes during evaporation in detail. Its structural status is subject to ongoing debate.

\subsection{Island formula and replica wormholes (2019--2022)}

Almheiri, Engelhardt, Marolf, and Maxfield [8] and subsequent work [9] developed the \emph{island formula} for the entanglement entropy of Hawking radiation. The formula generalizes Ryu-Takayanagi:
\[
S_{\mathrm{rad}}(t) = \min\left[ \mathrm{Ext}\left( \frac{A(\partial I)}{4 G_N} + S_{\mathrm{matter}}(R \cup I) \right) \right]
\]
where $I$ is the \emph{island} --- a region inside the BH that contributes to radiation entanglement entropy when extremized over candidate island geometries. The replica-wormhole construction provides a path-integral derivation: at the Page time, replica spacetimes connecting different copies of the BH dominate, and the resulting entanglement entropy follows the Page curve.

The structural assumptions on which the island formula depends:

\textbf{(A13) Replica-trick path integrals} as physically meaningful.

\textbf{(A14) Wormhole geometries} in the gravitational path integral, with non-trivial topology connecting BH copies.

\textbf{(A15) Quantum extremal surface} prescription (generalization of Ryu-Takayanagi-Hubeny-Rangamani-Takayanagi).

\textbf{(A16) Geometric horizon} as a boundary at which the entanglement-entropy formula evaluates.

The island formula is the most sophisticated current proposal and has reproduced the Page curve in specific calculations [9]. Its structural status is high; its dependence on path-integral wormholes and the geometric-horizon assumption (A16) places specific commitments on the fundamental ontology.

\subsection{The shared structural commitments}

Across the major proposals, the following commitments recur:

\begin{center}
\small
\begin{tabular}{@{}p{3.4cm}p{1.8cm}p{1.4cm}p{1.9cm}p{1.6cm}p{1.9cm}@{}}
\toprule
Assumption & Info paradox & Firewall & ER=EPR & Soft hair & Island formula \\
\midrule
Global unitarity & A1 & A1 & (assumed) & (assumed) & (assumed) \\
Geometric horizon as primitive & A2 & (implicit) & A8 (wormhole) & A12 & A16 \\
Monogamy at boundary & (implicit) & A5 & (resolved via A8) & (implicit) & (implicit) \\
Global Cauchy data & (implicit) & (implicit) & (implicit) & (implicit) & (implicit) \\
Smooth-horizon idealization & A4 & A6 & (resolved via A8) & A12 & A16 \\
Wormhole topology & --- & --- & A7 & --- & A14 \\
\bottomrule
\end{tabular}
\end{center}

The ED framework, as we develop in \S\S3--8, does not impose these commitments at the substrate level. Where they appear in our derivations, they appear as \emph{coarse-grained continuum-level features} emergent from substrate dynamics --- not as primitive structural commitments.

\section{Event Density Substrate Ontology (Compressed)}

For full development of the ED ontology, see [12, 13]. We summarize the load-bearing pieces here.

ED takes the universe to consist of discrete \emph{micro-events} (P01) linked by \emph{participation} relations into \emph{chains} (P02). Bandwidth (P04) is the graded measure of participation, additive for independent contributions. Polarity (P09) supplies a $U(1)$-valued phase relation. Continuous time (P13) governs evolution between commitments. Commitment irreversibility (P11) makes participation events permanent at the substrate level.

The framework's load-bearing kernels are:

\textbf{V1 finite-width vacuum kernel} (Theorem 18 [14]): mediates cross-chain correlations forward-in-time-only, with finite spatial width $\sim \ell_P$.

\textbf{V5 finite-memory kernel}: mediates substrate cross-chain memory effects, with finite temporal width $\tau_{V5}$.

The Diffusion Coarse-Graining Theorem (DCGT) [15] provides the substrate-to-continuum bridge via hydrodynamic-window scale separation $\ell_P \ll R_{\mathrm{cg}} \ll L_{\mathrm{flow}}$ and multi-scale expansion in $\ell_P/R_{\mathrm{cg}}$.

Theorem 19 (Newton-recovery) [16] identifies the substrate length scale $\ell_P$ as the Planck length via $G = c^3 \ell_P^2 / \hbar$.

Theorem 17 (gauge-field-as-rule-type) [17] identifies gauge fields as the substrate's rule-type connection, providing the structural framework for charged-particle Hawking emission.

Two interpretation-level structures are load-bearing:

\textbf{ED-I-06 (No fundamental fields)} [18]: the ED ontology does not posit fundamental field-theoretic objects. Fields appear as continuum-level coarse-grained features of substrate dynamics, not as primitive ontology.

\textbf{Substrate locality}: cross-chain interactions occur via mediating substrate structure (V1, V5, channels), not via instantaneous non-local action.

The framework's substrate ontology is therefore \emph{ontologically minimal}: discrete events, finite participation, irreversible commitment, finite-width kernels, locality. From this minimal substrate, continuum-level structures (spacetime, gauge fields, gravity, quantum mechanics) emerge as coarse-grained leading-order content of substrate dynamics. The standard structural commitments (A1--A16) are not imposed at the substrate level; they appear, in modified form, as continuum-level features when DCGT coarse-graining produces them.

\section{The Decoupling Surface as Substrate-Level Horizon}

In the ED framework, what classical general relativity calls a \emph{black-hole horizon} is, at the substrate level, a \emph{saturated decoupling surface}: a locus where the gradient sparsity $\sigma$ has reached its substrate-saturation value $\sigma_{\mathrm{sat}}$, with the cross-chain bandwidth across the surface exponentially suppressed below hydrodynamic-window resolution.

The mechanism: cross-chain bandwidth between substrate regions takes the form
\[
\Gamma_{\mathrm{cross}}(\mathbf{x}_1, \mathbf{x}_2) \sim \exp\left[-\alpha \int_{\mathrm{path}} \sigma(\mathbf{x})\, d\ell\right]
\]
where $\sigma(\mathbf{x}) = |\nabla\rho|\, \ell_P^2 / \rho_{\mathrm{local}}$ is the substrate gradient sparsity and $\alpha$ is a substrate-determined dimensionless prefactor. At the saturated surface, $\sigma \to \sigma_{\mathrm{sat}}$ rapidly along radial paths, producing exponential suppression of $\Gamma_{\mathrm{cross}}$.

When the suppression exceeds DCGT resolution --- formally, $|\nabla\rho|\, \ell_P^2 / \rho_{\mathrm{local}} \gtrsim \log(R_{\mathrm{cg}}/\ell_P)$ --- the substrate-to-continuum coarse-graining cannot resolve any nonzero cross-bandwidth. The decoupling surface is the substrate-level locus where this condition is met.

Several structural features of this construction are worth emphasizing:

\textbf{Statistical, not geometric.} The decoupling surface is a feature of the substrate's correlational structure, not of any underlying spacetime geometry. The surface emerges as a statistical feature of bandwidth-collapse, not as a primitive geometric boundary. The smooth-horizon idealization that the standard paradoxes (A2, A6, A12, A16) require is not present at the substrate level.

\textbf{Universal across horizon types.} The same substrate mechanism produces apparently distinct horizon types: BH horizons (gradient from localized mass), Rindler horizons (gradient from acceleration), cosmological horizons (gradient from cosmic expansion at scale $R_H = c/H_0$), and analog-gravity horizons (gradient in BEC, acoustic, or photonic analog systems). The unification is forced --- same mechanism, different gradient sources [10, 11].

\textbf{Substrate-level surface gravity.} The substrate analog of surface gravity is
\[
\kappa_{\mathrm{ED}} = \alpha \cdot (\nabla\sigma)|_{\mathrm{surf}}
\]
which identifies with the standard $\kappa$ at leading-order DCGT coarse-graining. For Schwarzschild: $\kappa_{\mathrm{ED}} \to \kappa = c^4/(4 G M)$.

\textbf{Geometric horizon as continuum-emergent shadow.} The standard event horizon of GR appears as the level set of the acoustic-metric null-cone tilting at the decoupling surface. It is the continuum-level reading of a substrate-level statistical feature. This dissolves the \emph{teleological} aspect of standard event-horizon definitions (which require global Cauchy data to specify) and replaces it with a local substrate-level criterion.

The decoupling surface is the substrate-level structure on which Arc Hawking's machinery operates.

\section{The Substrate-Unruh Derivation of Hawking Radiation}

The standard semiclassical derivation of Hawking radiation runs through Bogoliubov transformations between past and future modes of a quantum field on Schwarzschild spacetime [1]. Arc Hawking [10] develops a structurally parallel derivation using substrate primitives.

\subsection{The substrate-Unruh argument}

Near the saturated decoupling surface, introduce substrate coordinates $(\rho, t)$ with $r = \rho^2 \kappa_{\mathrm{ED}} / 2$. Substrate observers at fixed $\rho$ are accelerated relative to the substrate's asymptotic frame, with proper substrate-acceleration $a(\rho) = 1/\rho$.

The substrate-level Unruh effect: a substrate observer at fixed $\rho$ near the saturated surface sees the V5 vacuum as a thermal state at temperature
\[
T_{\mathrm{local}} = \frac{a}{2\pi} = \frac{1}{2\pi\rho}.
\]

This is structurally parallel to the standard Unruh effect [19], with the substrate's V5 vacuum playing the role of the standard Wightman vacuum.

\subsection{Imaginary-time periodicity}

Under analytic continuation $t \to -i\tau$, the substrate's local geometry near the saturated surface becomes
\[
ds^2_{\mathrm{substrate}} \approx \rho^2 \kappa_{\mathrm{ED}}^2 \, d\tau^2 + d\rho^2 + \cdots
\]

This is the substrate analog of flat 2D Euclidean space in polar coordinates. Avoiding a substrate-level conical singularity at $\rho = 0$ requires the angular coordinate $\kappa_{\mathrm{ED}}\tau$ to have period $2\pi$:
\[
\beta = 2\pi/\kappa_{\mathrm{ED}}.
\]

The substrate-level argument for this periodicity: V5 cross-chain correlations near the saturated surface must form a self-consistent substrate object. A conical singularity in the analytically-continued substrate would correspond to a discontinuity in V5 correlations as the imaginary substrate-time circle is traversed, which substrate locality (P02) and continuity-under-DCGT-coarse-graining forbid. The periodicity $\beta = 2\pi/\kappa_{\mathrm{ED}}$ is therefore \emph{forced} by substrate-vacuum self-consistency, in parallel with the standard no-conical-singularity argument in QFT-in-curved-spacetime [20].

\subsection{KMS condition and Planck distribution}

A correlation function periodic in imaginary time with period $\beta$ satisfies the KMS condition [21], mathematically equivalent to thermal correlations at temperature $T = 1/\beta$. The substrate's V5 correlation function across the saturated surface therefore satisfies the KMS condition at temperature
\[
T = \frac{\kappa_{\mathrm{ED}}}{2\pi}.
\]

The spectral content is the Planck distribution:
\[
N(\omega) = \frac{1}{e^{\beta\omega} - 1}, \qquad \omega > 0.
\]

\subsection{DCGT identification with Hawking temperature}

At leading-order DCGT coarse-graining, $\kappa_{\mathrm{ED}} \to \kappa$ (the standard surface gravity). For Schwarzschild, $\kappa = c^4/(4 G M)$, recovering the Hawking temperature
\[
T_H = \frac{\kappa}{2\pi} = \frac{\hbar c^3}{8\pi G M k_B}.
\]

The framework reproduces the Hawking temperature exactly at leading order via the substrate-Unruh argument plus DCGT identification.

\section{The V5 Cutoff and Trans-Planckian Resolution}

The standard semiclassical Hawking calculation is structurally challenged by the \emph{trans-Planckian problem}: modes observed at moderate frequencies at infinity arose from arbitrarily-blueshifted Planck-scale modes near the horizon, where standard QFT is structurally questionable [22]. The standard derivation assumes ordinary QFT applies at all scales --- empirically supported by the calculation's success but theoretically uncomfortable.

ED resolves this at the substrate level. The V5 finite-memory kernel has characteristic time $\tau_{V5} = \ell_P/c$ at the gravitational scale, identified via Theorem 19 plus substrate dimensional analysis (Planck time is the unique substrate-built timescale at the gravitational scale). The kernel's frequency-domain magnitude:
\[
|\tilde V_5(\omega)|^2 = \frac{(\mathcal{V}_0 \tau_{V5})^2}{1 + (\omega\tau_{V5})^2}.
\]

For $\omega \ll \omega_c \equiv 1/\tau_{V5} = c/\ell_P$ (the Planck frequency), V5 correlations are coherent. For $\omega \gtrsim \omega_c$, V5 coherence breaks down --- the substrate cannot mediate cross-chain correlations at frequencies approaching the Planck scale.

The implications for Hawking radiation: substrate modes near the horizon at proper frequencies $\omega_{\mathrm{proper}} \gtrsim \omega_c$ do not maintain V5-coherent structure. The trans-Planckian content of standard semiclassical calculations is \emph{resolved at the substrate level} by V5's finite memory. The substrate does not support arbitrarily-blueshifted modes near the horizon.

The substrate-cutoff modulates the leading-order Planck distribution at first subleading order:
\[
N_{ED}(\omega) = \frac{N_H(\omega)}{1 + (\omega\tau_{V5})^2} \approx N_H(\omega) \cdot \left[1 - (\omega\tau_{V5})^2 + O((\omega\tau_{V5})^4)\right].
\]

For Schwarzschild stellar mass, the correction at the spectrum's peak ($\omega \sim T_H$) is $(\omega\tau_{V5})^2 \sim (\ell_P/M)^2 \sim 10^{-76}$ --- entirely invisible. For primordial BHs in their final stages of evaporation, $(\ell_P/M)^2$ approaches order unity, and the correction becomes significant.

This is one of the framework's most theoretically significant contributions: a \emph{natural UV regulation} that resolves a structural problem of standard semiclassical Hawking without requiring phenomenological cutoffs or ad-hoc regulators. The substrate's V5 finite memory is the regulator.

\section{Greybody Factors, Page Rate, and Page Curve}

Arc Hawking [10] develops the substrate-derived greybody factors, integrated emission rate, and Page curve. We summarize the key results.

\subsection{Greybody factors}

The radial substrate wave equation:
\[
\frac{d^2 R_{\omega\ell}}{dr_*^2} + [\omega^2 - V_\ell^{(\mathrm{ED})}(r)] R_{\omega\ell} = 0
\]

with substrate-effective potential
\[
V_\ell^{(\mathrm{ED})}(r) = f_\sigma(r) \left[\frac{\ell(\ell+1)}{r^2} + \frac{1}{r}\frac{df_\sigma}{dr}\right]
\]

where $f_\sigma(r) = \exp[-\alpha \int_{r_h}^r \sigma(r')\, dr'/\ell_P^2]$ is the substrate state-dependent factor. At leading-order DCGT, $f_\sigma(r) \to 1 - 2M/r$, recovering the standard Regge-Wheeler effective potential [23].

The substrate transmission coefficient $\mathcal{T}_\ell^{(\mathrm{ED})}(\omega) = |T_\ell^{(\mathrm{ED})}(\omega)|^2$ identifies with the standard Regge-Wheeler / Teukolsky greybody factors $\mathcal{T}_\ell^{(\mathrm{GR})}(\omega)$ at leading order, with first-subleading-order corrections at $(\ell_P/M)^2$ from V5 cutoff and motif-alphabet effects.

\subsection{Page evaporation rate}

Integration of the leading-order spectrum and greybody factors over modes reproduces the standard Page rate [24]:
\[
\frac{dM}{dt}\bigg|_{\mathrm{leading}} = -\frac{\alpha_{\mathrm{Page}}}{M^2}.
\]

First-subleading-order corrections give:
\[
\frac{dM}{dt}\bigg|_{ED} = -\frac{\alpha_{\mathrm{Page}}}{M^2} \left[1 - K\left(\frac{\ell_P}{M}\right)^2 + O\left(\frac{\ell_P}{M}\right)^4\right]
\]

where $K = -c_{V5}^{\mathrm{int}} - c_g^{\mathrm{int}} \log g$ is inherited from substrate-microscopic details (V5 kernel parameters and BH-5 motif alphabet $g$ [25]).

\subsection{Page curve via bandwidth-budget min-bound}

Arc E (entanglement) [26] established that bipartite entanglement is bounded by the substrate's cross-chain bandwidth budget. For the Hawking system A (outgoing radiation) and B (interior modes), this becomes
\[
S_{AB}(t) \leq \min\left[S_{\mathrm{BH}}(t), S_{\mathrm{radiation}}(t)\right].
\]

This min-bound, combined with BH-5's area-law entropy $S_{\mathrm{BH}}(t) = (\log g)/(4\ell_P^2) \cdot A(t)$ and the cumulative radiation entropy $S_{\mathrm{radiation}}(t) \sim S_{\mathrm{BH,0}}\cdot t/\tau_{\mathrm{BH}}$, produces the Page curve [27]:

\begin{itemize}[noitemsep]
  \item Linear rise from 0 to $S_{\mathrm{max}} = S_{\mathrm{BH,0}}/2$ at $t_{\mathrm{Page}} \approx 0.54 \tau_{\mathrm{BH}}$
  \item Power-law fall from $S_{\mathrm{max}}$ toward 0 (or $S_{\mathrm{remnant}}$ in the remnant scenario) at $\tau_{\mathrm{BH}}$
  \item Cusp at $t_{\mathrm{Page}}$ where the min-bound switches argument
\end{itemize}

The substrate-level Page curve is reproduced exactly at leading-order DCGT coarse-graining.

\subsection{First-subleading-order corrections to Page time}

V5 cutoff and motif-alphabet effects modify Page time and maximum entropy:
\[
t_{\mathrm{Page}}^{(ED)} = 0.54\, \tau_{\mathrm{BH}} \cdot \left[1 + c_t \left(\frac{\ell_P}{M_0}\right)^2 + O\left(\frac{\ell_P}{M_0}\right)^4\right]
\]

with $c_t$ inherited from substrate-microscopic details. For stellar-mass BHs, the correction is $\sim 10^{-76}$; for primordial BHs in final stages, it becomes significant.

\section{The Bandwidth-Budget Mechanism and Cross-Domain Unification}

The bandwidth-budget mechanism is one of ED's most striking cross-domain unifications. The same substrate machinery produces:

\begin{itemize}[noitemsep]
  \item \textbf{Bipartite monogamy of entanglement} at the qubit-pair scale (Arc E, Memo 4 [26]).
  \item \textbf{Entanglement-straddling at black-hole horizons} at the macroscopic scale (BH-4 [11]).
  \item \textbf{Correlation-budget plateau} in high-multiplicity-redundancy quantum-computing architectures (Q-COMPUTE [28]).
  \item \textbf{Page-curve min-bound} at the BH-radiation bipartite scale (Arc Hawking, H-5 [10]).
\end{itemize}

Across these four contexts --- separated by many orders of magnitude in physical scale --- the same substrate mechanism operates: a finite local cross-chain bandwidth budget $\Gamma_{\mathrm{max}}$ at any substrate region, with the entanglement structure between regions bounded by the bandwidth available.

This is a \emph{substrate-level unification} of phenomena that standard physics treats as separate. The same machinery that enforces Coffman-Kundu-Wootters monogamy [29] for qubit triples enforces the Page-curve min-bound for BH evaporation. The substrate mechanism is the same; only the physical platform differs.

The cross-domain unification has a parallel in V5's role: the same V5 finite-memory kernel produces:

\begin{itemize}[noitemsep]
  \item \textbf{Maxwell viscoelastic memory} in soft matter (DCGT consequence [15]).
  \item \textbf{Hawking spectrum high-frequency cutoff} at the Planck scale (Arc Hawking H-4 [10]).
  \item \textbf{Entanglement-bandwidth modulation} in BH-radiation Page-curve evolution (H-5 [10]).
\end{itemize}

One substrate primitive, three physical applications, scale separation of $\sim 40$ orders of magnitude.

These cross-domain unifications constitute structural evidence for the framework's substrate ontology: phenomena that look unrelated in standard physics are unified by shared substrate machinery, with empirically falsifiable predictions in each sector.

\section{Information Architecture: Why ED Doesn't Generate the Paradoxes}

Combining the substrate-ontology of \S3, the decoupling-surface mechanism of \S4, the substrate-Unruh derivation of \S5, the V5 cutoff of \S6, and the Page-curve structure of \S\S7--8, we can now articulate why the standard paradoxes do not arise in the ED framework.

\subsection{Information vs. entanglement at substrate level}

The framework distinguishes two structurally different forms of correlation:

\textbf{Committed structure (information).} Patterns in the substrate that have undergone P11 commitment --- irreversibly recorded participation events. Committed structure transmission requires nonzero cross-chain bandwidth $\Gamma_{\mathrm{cross}}$ across the surface separating sender and receiver.

\textbf{Uncommitted structure (entanglement).} Joint participation-amplitude correlations between regions, set up by prior participation history but not requiring ongoing cross-bandwidth to persist.

At a saturated decoupling surface, $\Gamma_{\mathrm{cross}}$ is suppressed below DCGT resolution. Committed structure cannot cross. Entanglement (uncommitted structure) can straddle, because it doesn't require ongoing cross-bandwidth to persist.

This is the substrate translation of the standard quantum-information no-signaling theorem: a horizon blocks signaling but not pre-existing correlations.

\subsection{Why each standard assumption fails at substrate level}

\textbf{(A1) Global unitarity:} the framework is local and irreversible at substrate level via P11. Unitarity is a coarse-grained continuum property in the canonical hydrodynamic window for closed sub-systems, not a global substrate constraint. There is no global-unitarity requirement that must be reconciled with horizon-blocked information. \emph{Substrate-level unitarity} is over-determined by three independent locks (T18 forward-cone V1 kernel, P11 commitment irreversibility, ED-I-06 no-fundamental-fields) --- the same three locks that produce E-5's no-signaling theorem [26].

\textbf{(A2, A12, A16) Geometric horizon as primitive:} the horizon is a participation-bandwidth-collapse surface, not a smooth geometric boundary. The smooth-horizon idealization that the standard paradoxes require is not present at the substrate level.

\textbf{(A5) Monogamy at boundary:} the surface is not a knife-edge geometric locus, so there is nothing to enforce monogamy at. Monogamy of entanglement is itself substrate-derived (Arc E E-4 [26]) as a bandwidth-budget conservation law, not as a boundary condition imposed at a geometric horizon.

\textbf{(A4) Pure-thermal Hawking spectrum:} the radiation is \emph{thermal} in spectral form (Planck distribution at $T_H$), but it carries entanglement-correlation structure with the BH interior via the entanglement-straddling mechanism (BH-4). It is not ``purely random thermal'' --- the thermal spectrum has correlation structure encoded in the bipartite-entanglement bandwidth between A and B.

\textbf{(A6) Effective field theory smoothness:} the framework operates within the acoustic-metric scalar-tensor approximation [30]. EFT smoothness for an infalling observer is preserved at the coarse-grained continuum level; substrate-level dynamics remain finite throughout the saturated zone, with no firewall structure required.

\textbf{(A7, A8, A14) Wormhole topology:} ED-I-06 forbids fundamental fields and topological structures of this kind [18]. The substrate has no fundamental wormhole geometry. Cross-horizon entanglement structure is supplied by V5 cross-chain re-routing, not by wormhole identification.

\textbf{(A10, A11) Asymptotic symmetry algebra:} ED's substrate ontology has rule-type taxonomy [17] that plays a structurally analogous role to soft-hair charges (substrate-level encoding of microstate distinctions), but at the substrate level rather than as asymptotic boundary data.

\textbf{(A13, A15) Replica-trick path integrals and quantum extremal surfaces:} ED's bandwidth-budget min-bound (Arc E [26]) replaces the island-formula entanglement-entropy bookkeeping. The min-bound produces the Page curve directly from substrate primitives, without requiring path-integral wormholes or quantum-extremal-surface prescriptions.

\subsection{The no-paradox-generation framing}

The standard paradoxes arise from imposing structural assumptions (A1--A16) that ED's substrate ontology does not make. Without those assumptions, the paradoxes do not arise.

This is structurally distinct from other resolution programs. Firewalls \emph{resolve} the original information paradox by \emph{creating} the firewall paradox; ER=EPR \emph{resolves} the firewall by \emph{positing} wormholes; soft hair \emph{encodes} information at the boundary; the island formula \emph{generalizes} Ryu-Takayanagi to compute the Page curve. Each works within a paradox-generating framework and addresses subproblems within it.

ED operates outside the paradox-generating framework. The paradoxes are not solved within ED; they are not generated. The framework's contribution is \emph{no-paradox-generation}, not \emph{paradox-resolution}.

\section{Comparisons with Major Proposals}

\subsection{ED vs. Firewall}

\textbf{Firewall mechanism.} AMPS [5] argued that combining (A1) global unitarity, (A5) monogamy, and (A6) EFT smoothness for old BHs forces a high-energy wall at the horizon, contradicting the equivalence principle.

\textbf{ED mechanism.} No firewall arises in ED for two reasons:
\begin{itemize}[noitemsep]
  \item The horizon is not a geometric boundary at which monogamy can be enforced (A5 fails to apply).
  \item Monogamy is itself substrate-derived as a bandwidth-budget conservation (Arc E [26]); the bandwidth budget is finite at every substrate region but distributed across the saturated zone, not concentrated at a sharp boundary.
\end{itemize}

\textbf{Verdict.} ED resolves the firewall paradox by not generating it. The substrate-level bandwidth-budget structure handles the entanglement-monogamy bookkeeping without requiring boundary-localized monogamy enforcement.

\subsection{ED vs. ER=EPR}

\textbf{ER=EPR mechanism.} Maldacena-Susskind [6] proposed that entangled pairs are connected by Einstein-Rosen bridges; identification of interior and exterior modes via the wormhole geometry resolves the firewall by dissolving the apparent monogamy violation.

\textbf{ED mechanism.} The substrate has no fundamental wormhole topology (ED-I-06 [18]). Cross-horizon entanglement structure is supplied by V5 cross-chain re-routing around the saturated surface (BH-4 [11]). The mechanism is empirically indistinguishable from ER=EPR's predictions at leading order, but is structurally distinct: ED produces the cross-horizon entanglement structure via substrate-level V5 correlations rather than via geometric wormhole topology.

\textbf{Comparison.} ER=EPR posits geometric realization (A7, A8); ED produces the equivalent observable structure (entanglement straddling the surface) through substrate-level correlation re-routing without geometric realization. ED is structurally more economical: no fundamental wormhole geometry is required.

\textbf{Verdict.} ED reproduces ER=EPR's \emph{observable signature} (cross-horizon entanglement preservation) without committing to its \emph{ontology} (fundamental wormhole topology). Where ER=EPR requires (A7) wormhole topology, ED requires only V5 cross-chain re-routing --- already present in the substrate ontology.

\subsection{ED vs. Soft Hair}

\textbf{Soft hair mechanism.} HPS [7] proposed that BH microstates are distinguished by supertranslation and superrotation charges (soft hair) at the horizon and asymptotic infinity. Information falling into the BH is encoded in soft-hair imprints, allowing recovery via correlations with soft modes during evaporation.

\textbf{ED mechanism.} ED's rule-type taxonomy [17] plays a structurally analogous role: substrate-level microstate distinctions are encoded in rule-type connections (the substrate analog of gauge-field structure), with the BH-5 motif alphabet $g$ providing the substrate-counting structure for area-law entropy. Information falling into the BH is encoded in substrate-level participation-pattern correlations, with the entanglement-straddling mechanism transporting these correlations to the radiation field.

\textbf{Comparison.} Soft hair encodes information at \emph{asymptotic boundaries} (null infinity, horizon as boundary); ED encodes information in \emph{substrate-level participation-pattern structure} (rule-type connections, motif-alphabet, V5 cross-chain correlations). The structural commitments differ: soft hair requires (A10) asymptotic symmetry algebra and (A11) soft-mode encoding; ED requires only the substrate ontology.

\textbf{Verdict.} ED and soft hair are structurally compatible at leading order in the sense that both supply microstate-distinguishing structure at the boundary of the BH region. ED's substrate-level mechanism is independent of the asymptotic-symmetry framework; soft hair's mechanism does not require ED's substrate ontology. They could in principle be empirically distinguished by precision tests of microstate-counting at the asymptotic versus substrate levels.

\subsection{ED vs. Island Formula}

\textbf{Island formula mechanism.} AEMM [8] and replica-wormhole calculations [9] derive the Page curve via the entanglement-entropy formula
\[
S_{\mathrm{rad}}(t) = \min\left[ \mathrm{Ext}\left( \frac{A(\partial I)}{4 G_N} + S_{\mathrm{matter}}(R \cup I) \right) \right]
\]
where the \emph{island} $I$ is a region inside the BH that contributes to radiation entanglement entropy. The formula requires (A13) replica-trick path integrals, (A14) wormhole geometries, (A15) quantum extremal surface prescription, and (A16) geometric horizon as boundary.

\textbf{ED mechanism.} ED's bandwidth-budget min-bound (Arc E E-4 [26]) replaces the island-formula bookkeeping:
\[
S_{\mathrm{rad}}(t) \leq \min\left[S_{\mathrm{BH}}(t), S_{\mathrm{radiation}}(t)\right].
\]

This is structurally simpler than the island formula. It directly produces the Page curve from substrate-level bandwidth conservation without invoking replica-wormhole path integrals or quantum-extremal-surface prescriptions. The ``island'' content --- the contribution of interior modes to radiation entanglement entropy --- is supplied by the bandwidth-budget structure, which is already substrate-derived.

\textbf{Comparison.} The island formula is a sophisticated and successful framework, having reproduced the Page curve in specific calculations [9]. Its structural commitments are substantial: replica-trick gravitational path integrals plus wormhole topology plus the quantum-extremal-surface prescription. ED's bandwidth-budget approach makes none of these commitments and produces the Page curve from substrate-level primitives.

\textbf{Verdict.} ED and the island formula produce the same Page curve at leading order. The structural mechanisms differ: island formula uses geometric/topological path-integral machinery; ED uses substrate-level bandwidth-conservation. ED is more economical; the island formula is more developed within standard quantum-gravity language. The two could be empirically distinguished by precision tests of:
\begin{itemize}[noitemsep]
  \item The substrate-cutoff predictions (V5 high-frequency cutoff, motif-alphabet temperature shift) that ED predicts but the island formula does not naturally produce.
  \item The Planck-mass-remnant scenario (Scenario C in ED) that ED predicts conditionally but the island formula does not naturally produce.
\end{itemize}

\subsection{Comparison summary table}

\begin{center}
\scriptsize
\begin{tabular}{@{}p{2.4cm}p{1.9cm}p{1.7cm}p{1.9cm}p{1.6cm}p{1.7cm}p{2.0cm}@{}}
\toprule
Feature & Standard semiclassical & Firewall & ER=EPR & Soft hair & Island formula & ED \\
\midrule
Information loss & Predicted & Avoided & Avoided & Avoided & Avoided & Not generated \\
Horizon structure & Geometric boundary & Firewall & Wormhole identification & Soft-hair charges & QES at boundary & Substrate decoupling surface \\
Trans-Planckian regulation & Phenomenological & Same & Same & Same & Same & V5 substrate cutoff \\
Page curve mechanism & Not derived & Not derived & Not directly & Not directly & QES + replica wormholes & Bandwidth-budget min-bound \\
Information-recovery channel & Thermal correlations & Firewall radiation & Wormhole identification & Soft-hair correlations & Island contributions & V5 cross-chain re-routing \\
Smooth-horizon assumption & Required & Modified & Required (w/ wormhole) & Required & Required & Not required \\
Wormhole topology & Not required & Not required & Required & Not required & Required & Not present \\
\bottomrule
\end{tabular}
\end{center}

\section{Late-Time Scenario: Planck-Mass Remnant FORCED}

The ED framework's higher-order resummation of $(\ell_P/M)^2$ corrections to the Page rate (H-8 [10a]) produces a definite verdict on the late-time evaporation profile:

\begin{itemize}[noitemsep]
  \item \textbf{Scenario A (Full Recovery to $M = 0$):} \textbf{REFUTED} by substrate constraints. P4 bandwidth-finiteness + DCGT hydrodynamic-window closure + P11 commitment-irreversibility jointly forbid full evaporation to zero mass.
  \item \textbf{Scenario B (Modified Turnover, no remnant):} \textbf{INCOMPLETE.} V5-alone analysis produces Scenario-B-like slowing, but the full substrate inventory (V5 + DCGT + P4 + P11 + N1 + T19) produces Scenario C. V5 alone misses the DCGT hydrodynamic-window closure transition.
  \item \textbf{Scenario C (Stable Planck-mass remnant):} \textbf{FORCED at substrate-structural level.} The combination of V5 finite-memory cutoff + DCGT hydrodynamic-window closure + P4 bandwidth-finiteness + P11 commitment-irreversibility produces a stable substrate-determined endpoint at $M_* = c_*\ell_P$ with $c_*$ order-unity (INHERITED from substrate microscopic details). The substrate-saturated cluster at $M = M_*$ is permanent (P11-locked), DCGT-closed (substrate-to-continuum bridge fails at $L_{\mathrm{flow}} \sim \ell_P$), and P4-bounded (finite participation density).
\end{itemize}

The framework's late-time prediction is \textbf{definite}, not conditional. ED predicts stable Planck-mass remnants from PBH evaporation. Each remnant has:
\begin{itemize}[noitemsep]
  \item Mass $\sim 22\,\mu\mathrm{g}$ (order-unity coefficient INHERITED).
  \item Stored entropy $\sim O(\log g)$ bits per BH-5's motif-alphabet structure.
  \item Substrate-bound information about parent BH initial state.
  \item Page curve asymptote at $S_{\mathrm{BH,0}} - S_{\mathrm{remnant}}$ rather than returning to zero.
\end{itemize}

\subsection*{Cosmological relic-matter abundance}

The cosmological abundance of these remnants is computed in H-9 [10b] as a \emph{structural cosmological prediction}. The present-day relic-matter fraction $\Omega_{\mathrm{relic}}$ depends on the empirical PBH formation history:

\begin{center}
\small
\begin{tabular}{@{}p{6.8cm}p{4.6cm}@{}}
\toprule
PBH formation scenario & $\Omega_{\mathrm{relic}}$ range \\
\midrule
Scale-invariant ($\alpha = 5/2$, Carr-Hawking standard) & $10^{-30}$ to $10^{-20}$ (negligible) \\
Critical collapse (Niemeyer-Jedamzik) & $10^{-41}$ to $10^{-25}$ (negligible) \\
Blue-tilted spectrum (moderate enhancement) & $10^{-15}$ to $10^{-5}$ (subdominant) \\
Inflationary-spike (low-mass concentration) & up to $10^{-3}$ (marginally compatible) \\
Inflationary-spike (extreme) & $\gtrsim 10^{-3}$ (largely excluded) \\
\bottomrule
\end{tabular}
\end{center}

For most observationally compatible scenarios, $\Omega_{\mathrm{relic}} \ll 1$ and the remnant population is cosmologically subdominant.

\subsection*{Critical framing note: not a dark-matter candidate}

ED's substrate-gravity content (separate galactic\_dynamics walkthrough [10c]) explains the empirical phenomenology of galactic dynamics --- flat rotation curves, slope-4 BTFR, transition acceleration $a_0 = cH_0/(2\pi)$, radial-acceleration relation --- \textbf{without invoking dark matter}. The framework does not require, and does not posit, a dark-matter particle population to explain galactic dynamics. What $\Lambda$CDM attributes to dark-matter halos, ED attributes to substrate-level dipole-mode projection plus geometric-mean composition.

The Planck-mass remnant population is therefore \emph{not} a dark-matter candidate. It is a \emph{separate structural cosmological prediction} about cosmic relic-matter content, derived from the substrate-cutoff endpoint of PBH evaporation. The relic-matter component is observationally testable through CMB-distortion, gamma-ray-background, and structure-formation searches that constrain the parent PBH population and thereby the relic abundance.

\subsection*{Comparison with standard proposals}

The standard proposals (firewall, ER=EPR, soft hair, island formula) do not naturally produce a Planck-mass-remnant scenario. The standard semiclassical Hawking calculation predicts complete evaporation in finite time. ED's regulated-completion framework produces Scenario C as a substrate-FORCED structural feature --- a definite prediction rather than a phenomenological possibility.

\section{Falsifiable Predictions}

The framework produces falsifiable predictions across multiple empirical platforms.

\subsection{Analog Hawking experiments (most accessible)}

In analog-gravity systems (BEC, acoustic, photonic), the substrate-cutoff scale corresponds to the analog system's microscopic correlation length. ED predicts:
\begin{itemize}[noitemsep]
  \item High-frequency cutoff at $\omega \sim 1/\tau_{\mathrm{analog}}$ in the analog Hawking spectrum.
  \item Cutoff form: $1/(1 + (\omega\tau_{\mathrm{analog}})^2)$ --- V5-derived structure.
  \item Spectrum-integrated emission rate corrected at $(T_H^{\mathrm{analog}}/\omega_c^{\mathrm{analog}})^2$.
\end{itemize}

Existing experiments [32] confirm spectral form at moderate frequencies; precision tests at the cutoff scale are within reach of current-generation analog Hawking experimental programs.

\subsection{Primordial BH late-stage evaporation (mid-range)}

For PBHs in their final stages of evaporation:
\begin{itemize}[noitemsep]
  \item Spectral cutoff at $\omega \sim \omega_c = c/\ell_P$ (Planck frequency).
  \item Modified evaporation profile: slowed in the final $\sim 0.1$ s at $M \sim M_P$.
  \item Possible Planck-mass remnant (Scenario C).
\end{itemize}

No PBH evaporations have been detected to date [33]. Continued absence at improving sensitivity (Fermi, HESS, CTA) constrains the framework's prediction. Detection would test the substrate-cutoff prediction directly.

\subsection{High-energy photon dispersion (mid-range)}

The substrate-cutoff scale $\omega_c = c/\ell_P$ produces dispersion in vacuum-substrate propagation of high-energy astrophysical photons:
\[
\Delta v / c \sim (\omega/\omega_c)^2.
\]

Public LIGO/Virgo + Fermi-LAT data on high-energy photons from gamma-ray bursts and gravitational-wave events permit precision time-of-flight measurements. The framework's prediction connects with E2 (GRB photon-timing retrodiction) in the program's Investigation Priority List [34].

\subsection{Relic-matter signature searches (long-range)}

The Planck-mass remnant population (FORCED via H-8) constitutes a structural relic-matter component of the cosmic energy budget. It is \emph{not} a dark-matter candidate --- ED's substrate-gravity already explains galactic dynamics without dark matter (galactic\_dynamics walkthrough [10c]).

The relic-matter abundance $\Omega_{\mathrm{relic}}$ (computed in H-9 [10b]) is testable through specific cosmological signatures:
\begin{itemize}[noitemsep]
  \item Microlensing: not feasible (Planck-mass too small for current sensitivity).
  \item Cosmic-ray collisions involving Planck-mass relics: extreme-energy detector signatures (current sensitivity insufficient).
  \item Structure-formation effects: weak gravitational signatures (dilute population in standard scenarios).
  \item BBN, CMB-distortion, and gamma-ray-background constraints on the parent PBH population indirectly bound the relic abundance.
\end{itemize}

Continued absence of Planck-mass relic-matter signatures at improving experimental sensitivity continues to constrain $\Omega_{\mathrm{relic}}$ from above. This constrains the \emph{relic-matter component} of the cosmic energy budget, not dark matter (which the framework does not posit). For most observationally compatible PBH formation scenarios, $\Omega_{\mathrm{relic}} \ll 1$ and the remnant population is cosmologically subdominant.

\subsection{Falsifiability summary}

\begin{center}
\small
\begin{tabular}{@{}p{4.4cm}p{3.6cm}p{2.6cm}p{2.4cm}@{}}
\toprule
Prediction & Platform & Accessibility & Status \\
\midrule
V5 cutoff form $1/(1+(\omega\tau)^2)$ & Analog Hawking (BEC, acoustic) & Current generation & Testable \\
Planck-scale spectral cutoff in PBH evaporation & Gamma-ray observatories & Detection-pending & Falsifiable \\
High-energy photon dispersion $(\omega/\omega_c)^2$ & LIGO/Virgo + Fermi-LAT & Current data & Closeable \\
Planck-mass remnant signatures & Direct/indirect searches & Long-term & Constraints accumulating \\
\bottomrule
\end{tabular}
\end{center}

\section{The Regulated-Completion Verdict}

The framework's relationship to standard semiclassical Hawking is best characterized as \emph{regulated completion}. The structural classification:

\textbf{Strict extension at observable scales.} ED reproduces every standard semiclassical Hawking ingredient exactly at leading-order DCGT coarse-graining. At observable BH scales (stellar mass, $M \gtrsim M_\odot$, $(\ell_P/M)^2 \sim 10^{-76}$), the framework is empirically indistinguishable from semiclassical.

\textbf{Substrate-level UV regulation.} The V5 finite-memory cutoff at $\tau_{V5} = \ell_P/c$ provides natural UV regulation of the standard semiclassical calculation, resolving the trans-Planckian problem at the substrate level rather than via phenomenological cutoffs.

\textbf{Modified-theory content at extreme scales.} At extreme scales (Planck-mass-approaching BHs), the framework predicts qualitatively distinct content: V5 cutoff dominance in the spectrum, \textbf{Scenario C Planck-mass remnant FORCED via H-8 substrate-cutoff resummation}, modified late-time evaporation profile. The remnant population constitutes a structural relic-matter component of the cosmic energy budget (not a dark-matter candidate; ED's substrate-gravity already handles galactic dynamics without DM).

\textbf{No replacement of semiclassical Hawking.} The framework includes all of semiclassical at leading order; it does not replace it. It grounds it.

The verdict: ED is a \emph{regulated completion} of semiclassical Hawking --- strict extension at observable scales, substrate-level UV regulation, possibly modified theory at extreme scales (Scenario C).

This positions the framework constructively relative to the active BH-information literature. ED does not compete with standard semiclassical Hawking; it provides the substrate-level account that standard semiclassical lacks. ED does not compete with the major resolution proposals (firewall, ER=EPR, soft hair, island formula); it operates at a different structural level --- substrate ontology versus continuum-field-theoretic ontology.

The framework's principal claim is that the standard paradoxes are not generated when the underlying ontology is substrate-discrete with finite participation, irreversible commitment, finite-width kernels, and locality. Without the smooth-horizon idealization on which the paradoxes structurally depend, the paradoxes simply do not arise.

\section{Open Questions and Future Work}

Several structural questions remain open:

\textbf{Higher-order resummation for the late-time scenario.} Whether Scenario A, B, or C realizes requires resummation of $(\ell_P/M)^2$ corrections at all orders. This is the most cosmologically significant deferred work.

\textbf{Substrate-microscopic derivation of the motif alphabet $g$.} The Bekenstein-Hawking $\log g \approx 1/4$ coefficient is currently inherited rather than derived. Closed-form V1-kernel motif counting would resolve this.

\textbf{Extension to charged BHs (Reissner-Nordstr\"om substrate analog).} Requires T17 (gauge-field-as-rule-type) inheritance for charged-particle Hawking emission and modified substrate-state-dependent factor at the charged horizon.

\textbf{Extension to rotating BHs (Kerr substrate analog).} Requires generalization of the saturated-surface substrate state. Couples with O1 (superradiance amplitude derivation) and O3 (full Kerr interior audit) in the program's Investigation Priority List [34].

\textbf{Beyond the acoustic-metric scalar-tensor approximation.} ED operates within this approximation [30]; full GR remains speculative. Strong-curvature regimes beyond acoustic-metric extension are out of current scope.

\textbf{Empirical engagement with analog Hawking experimental programs.} Current-generation analog systems can test the V5 cutoff form prediction. Experimental design memos coordinating with active BEC and acoustic Hawking experimental programs would advance falsifiability.

These open items do not undermine the framework's substrate-level resolution of the BH information question, but they identify directions for continued development.

\section{Conclusions}

The black-hole information paradox has resisted resolution within standard semiclassical and quantum-gravitational frameworks for nearly five decades. The major proposals --- firewall, ER=EPR, soft hair, island formula --- each address subproblems within paradox-generating frameworks while accepting structural commitments (global unitarity, geometric horizon, monogamy at boundary, replica-wormhole path integrals) that constrain the available solution space.

The Event Density framework operates at a different structural level. The substrate ontology of discrete micro-events on chains, with finite participation bandwidth, irreversible commitment-events, and finite-width kernels, does not impose the structural commitments on which the standard paradoxes depend. The framework does not resolve the paradoxes within a paradox-generating structure; it does not generate them.

Specifically:
\begin{itemize}[noitemsep]
  \item The horizon is a substrate-level statistical feature of bandwidth collapse, not a smooth geometric primitive.
  \item Information cannot cross the horizon by participation-bandwidth prohibition, not causal prohibition.
  \item Entanglement straddles the horizon via V5 cross-chain re-routing.
  \item Evaporation is participation re-routing through pair-creation-class events at the saturated surface.
  \item The Page curve is reproduced via Arc E's bandwidth-budget min-bound, without requiring replica-wormhole path integrals.
  \item The trans-Planckian problem is regulated at the substrate level via V5 finite memory at $\tau_{V5} = \ell_P/c$.
  \item Substrate-level unitarity is over-determined by three independent locks (T18 + P11 + ED-I-06).
\end{itemize}

The framework reproduces standard semiclassical Hawking exactly at leading-order DCGT coarse-graining. At first subleading order, it produces FORM-FORCED corrections at $(\ell_P/M)^2$ from V5 cutoff and motif-alphabet effects, with COEFFICIENTS-INHERITED from substrate-microscopic details. At extreme scales (Planck-mass-approaching BHs), the H-8 higher-order resummation analysis establishes that \textbf{Scenario C --- stable Planck-mass remnant --- is FORCED at substrate-structural level}. The H-9 relic-abundance calculation gives $\Omega_{\mathrm{relic}}$ scenario-dependent on PBH formation history, ranging from negligible ($10^{-30}$) for standard scale-invariant spectra to $\sim 10^{-3}$ for inflationary-spike scenarios at the boundary of observational viability. \textbf{The framework does not invoke dark matter at any stage:} substrate-gravity (separate galactic\_dynamics walkthrough) explains galactic dynamics, and the Planck-mass remnant population is a \emph{separate} structural relic-matter component of the cosmic energy budget.

The framework's relationship to standard semiclassical Hawking is \emph{regulated completion}: strict extension at observable scales, substrate-level UV regulation, possibly modified theory at extreme scales. The framework does not replace semiclassical Hawking; it grounds it.

The empirical predictions distinguishing ED from the major proposals are accessible:
\begin{itemize}[noitemsep]
  \item The V5 cutoff form $1/(1+(\omega\tau)^2)$ in analog Hawking spectra (current generation).
  \item High-energy photon dispersion at the Planck scale (closeable with public data).
  \item Modified primordial-BH evaporation profile in final stages (detection-pending).
  \item Possible Planck-mass remnant signatures (long-term).
\end{itemize}

The framework's substrate-level resolution of the black-hole information question --- based on no-paradox-generation rather than paradox-resolution-within-a-paradox-generating-framework --- represents a structurally distinct contribution to the active literature. The phrasing that captures the framework's stance: \emph{black holes do not have paradoxes because black holes do not have smooth edges, because the substrate is discrete}.

\section*{Section-by-Section Summary}

\begin{center}
\small
\begin{tabular}{@{}cp{4.6cm}p{8.2cm}@{}}
\toprule
\S & Title & Core content \\
\midrule
1 & Introduction & Paradox landscape; ED's distinctive claim of no-paradox-generation \\
2 & Standard paradoxes and assumptions & A1--A16 across major proposals; shared structural commitments \\
3 & ED substrate ontology & P01--P13, V1/V5 kernels, DCGT, T17/T18/T19, ED-I-06 \\
4 & Decoupling surface as substrate horizon & Bandwidth collapse; geometric horizon as continuum shadow \\
5 & Substrate-Unruh derivation & Imaginary-time periodicity $\to$ KMS $\to$ Planck $\to$ DCGT $\to$ $T_H$ \\
6 & V5 cutoff and trans-Planckian resolution & $\tau_{V5} = \ell_P/c$; substrate-level UV regulation \\
7 & Greybody, Page rate, Page curve & Regge-Wheeler at leading order; bandwidth-budget min-bound \\
8 & Bandwidth-budget mechanism & Cross-domain unification (qubits, BH, QC, soft matter) \\
9 & Information architecture & Why each paradox-generating assumption fails at substrate level \\
10 & Comparisons with major proposals & Firewall / ER=EPR / soft hair / island formula \\
11 & Late-time scenarios & A/B/C branching; Planck-mass remnant FORCED \\
12 & Falsifiable predictions & Analog Hawking, PBH evaporation, photon dispersion, relic-matter \\
13 & Regulated-completion verdict & Strict extension + UV regulation + conditional modified theory \\
14 & Open questions & Higher-order resummation, charged/rotating BHs, beyond acoustic-metric \\
15 & Conclusions & No-paradox-generation; structural distinctness from major proposals \\
\bottomrule
\end{tabular}
\end{center}

\section*{References}

\begin{enumerate}[leftmargin=*,itemsep=2pt]
  \item Hawking, S. W. ``Particle Creation by Black Holes.'' \emph{Communications in Mathematical Physics} \textbf{43}, 199--220 (1975).
  \item Bekenstein, J. D. ``Black Holes and Entropy.'' \emph{Physical Review D} \textbf{7}, 2333--2346 (1973).
  \item Hawking, S. W. ``Breakdown of Predictability in Gravitational Collapse.'' \emph{Physical Review D} \textbf{14}, 2460--2473 (1976).
  \item Susskind, L., Thorlacius, L., Uglum, J. ``The Stretched Horizon and Black Hole Complementarity.'' \emph{Physical Review D} \textbf{48}, 3743--3761 (1993).
  \item Almheiri, A., Marolf, D., Polchinski, J., Sully, J. ``Black Holes: Complementarity or Firewalls?'' \emph{Journal of High Energy Physics} \textbf{2013} (2), 062 (2013).
  \item Maldacena, J., Susskind, L. ``Cool Horizons for Entangled Black Holes.'' \emph{Fortschritte der Physik} \textbf{61}, 781--811 (2013).
  \item Hawking, S. W., Perry, M. J., Strominger, A. ``Soft Hair on Black Holes.'' \emph{Physical Review Letters} \textbf{116}, 231301 (2016).
  \item Almheiri, A., Engelhardt, N., Marolf, D., Maxfield, H. ``The Entropy of Bulk Quantum Fields and the Entanglement Wedge of an Evaporating Black Hole.'' \emph{Journal of High Energy Physics} \textbf{2019} (12), 063 (2019).
  \item Almheiri, A., Hartman, T., Maldacena, J., Shaghoulian, E., Tajdini, A. ``Replica Wormholes and the Entropy of Hawking Radiation.'' \emph{Journal of High Energy Physics} \textbf{2020} (5), 013 (2020).
  \item Proxmire, A. \emph{Arc Hawking: V5-Mediated Hawking Spectrum, Greybody Factors, Page Rate, Page Curve, and Substrate-Level Resolution of the Information Paradox.} May 2026.
  \item[{[10a]}] Proxmire, A. \emph{Arc Hawking H-8: Higher-Order Resummation and the Late-Time Evaporation Endpoint.} May 2026.
  \item[{[10b]}] Proxmire, A. \emph{Arc Hawking H-9: PBH Remnant Relic-Abundance.} May 2026.
  \item[{[10c]}] Proxmire, A. \emph{Walkthrough: From Primitives to Galactic Dynamics.} May 2026.
  \item[{[11]}] Proxmire, A. \emph{Arc BH (Black Hole Architecture).} April--May 2026.
  \item Proxmire, A. \emph{Event Density: One Substrate, Three Domains.} April 2026.
  \item Proxmire, A. \emph{Event Density Foundations: A Unified Substrate Architecture for Quantum, Fluid, Gauge, and Gravitational Dynamics.} April 2026.
  \item Proxmire, A. \emph{Theorem 18: V1 Kernel Retardation and the Kernel-Level Arrow of Time.} April 2026.
  \item Proxmire, A. \emph{The Diffusion Coarse-Graining Theorem.} April 2026.
  \item Proxmire, A. \emph{Theorem 19: Newton's Law from Substrate Holographic Counting.} April 2026.
  \item Proxmire, A. \emph{Theorem 17: Gauge-Field-as-Rule-Type.} April 2026.
  \item Proxmire, A. \emph{ED-I-06: Fields and Forces in Event Density.} February 2026.
  \item Unruh, W. G. ``Notes on Black-Hole Evaporation.'' \emph{Physical Review D} \textbf{14}, 870--892 (1976).
  \item Gibbons, G. W., Hawking, S. W. ``Action Integrals and Partition Functions in Quantum Gravity.'' \emph{Physical Review D} \textbf{15}, 2752--2756 (1977).
  \item Kubo, R. \emph{J. Phys. Soc. Japan} \textbf{12}, 570--586 (1957); Martin, P. C., Schwinger, J. \emph{Phys. Rev.} \textbf{115}, 1342--1373 (1959).
  \item Jacobson, T. ``Black-Hole Evaporation and Ultrashort Distances.'' \emph{Physical Review D} \textbf{44}, 1731--1739 (1991).
  \item Regge, T., Wheeler, J. A. ``Stability of a Schwarzschild Singularity.'' \emph{Physical Review} \textbf{108}, 1063--1069 (1957).
  \item Page, D. N. ``Particle Emission Rates from a Black Hole.'' \emph{Physical Review D} \textbf{13}, 198--206 (1976).
  \item Proxmire, A. \emph{Arc BH-5: Area-Law Entropy and the Motif Alphabet $g$.} April 2026.
  \item Proxmire, A. \emph{Arc E (Entanglement).} May 2026.
  \item Page, D. N. ``Information in Black Hole Radiation.'' \emph{Physical Review Letters} \textbf{71}, 3743--3746 (1993).
  \item Proxmire, A. \emph{Arc Q-COMPUTE: UR-1 Theorem and the Multiplicity-Cap Function.} May 2026.
  \item Coffman, V., Kundu, J., Wootters, W. K. ``Distributed Entanglement.'' \emph{Physical Review A} \textbf{61}, 052306 (2000).
  \item Proxmire, A. \emph{Arc ED-10: Curvature Emergence and the Acoustic-Metric Scalar-Tensor Approximation.} April 2026.
  \item Carr, B., Kuhnel, F., Sandstad, M. ``Primordial Black Holes as Dark Matter.'' \emph{Physical Review D} \textbf{94}, 083504 (2016).
  \item Steinhauer, J. \emph{Nature Physics} \textbf{12}, 959--965 (2016); Drori, J., et al. \emph{Physical Review Letters} \textbf{122}, 010404 (2019).
  \item Carr, B. J., Hawking, S. W. \emph{MNRAS} \textbf{168}, 399--415 (1974).
  \item Investigation Priority List, Event Density program archive, 2026.
\end{enumerate}

The full Event Density corpus is available at \url{https://github.com/allen-proxmire/event-density}.

\end{document}
